\documentclass[11pt,a4paper]{article}

\usepackage[utf8]{inputenc}
\usepackage[T1]{fontenc}
\usepackage{amsmath,amssymb,amsthm}
\usepackage{mathtools}
\usepackage{booktabs}
\usepackage{hyperref}
\usepackage[margin=1in]{geometry}

\newtheorem{theorem}{Theorem}[section]
\newtheorem{lemma}[theorem]{Lemma}
\newtheorem{proposition}[theorem]{Proposition}
\newtheorem{corollary}[theorem]{Corollary}
\theoremstyle{definition}
\newtheorem{definition}[theorem]{Definition}
\newtheorem{remark}[theorem]{Remark}

\DeclareMathOperator{\perm}{perm}
\DeclareMathOperator{\dc}{dc}
\DeclareMathOperator{\codim}{codim}
\DeclareMathOperator{\Sing}{Sing}
\DeclareMathOperator{\rank}{rank}
\DeclareMathOperator{\corank}{corank}
\DeclareMathOperator{\str}{str}
\DeclareMathOperator{\tr}{tr}
\DeclareMathOperator{\adj}{adj}
\DeclareMathOperator{\supp}{supp}
\DeclareMathOperator{\Span}{span}

\title{Word-Reversal Symmetry, Darboux Block Obstructions, and a General $\codim(\Sing) + 2$ Determinantal Complexity Bound for Degree-$n \ge 4$ Hypersurfaces\\{\large (with a Self-Contained Application to $\dc_{\mathbb{C}}(\perm_4) \ge 10$)}}
\author{BrientH}
\date{\today}

\begin{document}
\maketitle

\begin{abstract}
We develop an algebraic symmetry and obstruction framework explaining how origin-isotropy traps matrix operator chains in higher-order singular strata. First, in the free associative algebra $\mathcal{A} = \mathbb{C}\langle \mathbf{m}_1,\dots,\mathbf{m}_s \rangle$ equipped with the word-reversal anti-involution $\rho(\mathbf{m}_{i_1}\cdots\mathbf{m}_{i_k}) = \mathbf{m}_{i_k}\cdots\mathbf{m}_{i_1}$, we prove that any representation $\Phi : \mathcal{A} \to \mathbb{C}^{d_v\times d_v}$ satisfying the two initial skew-symmetry conditions $K^T = -K$ and $(\Phi(\mathbf{m}_i)K)^T = -\Phi(\mathbf{m}_i)K$ satisfies the universal identity $(\Phi(P)K)^T = -\Phi(\rho(P))K$ for all $P \in \mathcal{A}$. We establish necessary and sufficient conditions for quadratic annihilation ($y^T \Phi(P)K y \equiv 0 \iff \Phi(P^-)K = 0$), derive exact dimension formulas $\dim \mathcal{A}_k^\pm = (s^k \pm s^{\lceil k/2 \rceil})/2$, and characterize the admissible operator space $\mathcal{L}_K$ in a Darboux basis of $K$ ($\rank(K) = 2r \le d_v$).

Second, we apply this characterization to Algebraic Branching Programs (ABPs) and linear time-invariant (LTI) control systems: we show that two vanishing initial Markov parameters force the entire infinite Hankel matrix $H_\infty$ to collapse to zero, prove a width lower bound $w \ge \codim\Sing(V(F)) + 1$ for time-homogeneous and palindromic ABPs, and prove that even for completely non-homogeneous ABPs with arbitrary distinct layer matrices $A^{(1)} \neq \cdots \neq A^{(n-2)}$, the Hessian rank on $W^\perp$ is trapped by $\rank\mathrm{Hess}(F)(\xi) \le 2r = \rank(K) \le w - d_z$.

Third, we prove a general unconditional determinantal complexity theorem extending the Alper--Bogart--Velasco (2017) bound ($\dc_{\mathbb{C}}(F) \ge \codim\Sing(V(F)) + 1$): for \emph{every} homogeneous polynomial $F \in \mathbb{C}[x_1,\dots,x_N]_n$ of degree $n \ge 4$ satisfying $\codim_{\mathbb{C}^N}(\Sing(V(F))) \ge 5$ and $\codim_{\mathbb{C}^N}(\Sing_2(V(F))) > \codim_{\mathbb{C}^N}(\Sing(V(F)))$, one has unconditionally
\[
  \dc_{\mathbb{C}}(F) \;\ge\; \codim_{\mathbb{C}^N}\big(\Sing(V(F))\big) + 2.
\]
As a self-contained concrete application (combining an exact degree-$9$ finite-field Macaulay certificate $\codim\Sing(V(\perm_4)) = 8$ with a hand proof of $\codim\Sing_2(V(\perm_4)) = 12$), this immediately yields $\dc_{\mathbb{C}}(\perm_4) \ge 10$. (We note that a November 2026 preprint by Sheshadri, \texttt{arXiv:2611.06658}, not yet peer-reviewed, describes $10$ in its abstract as the best known lower bound without specifying the field and claims $\dc_{\mathbb{C}}(\perm_4) = 15$; our main contribution is the general degree-$n \ge 4$ hypersurface theorem and the free-algebra/ABP Darboux obstruction framework.)
\end{abstract}

\section{Introduction}

Let $X = (x_{i,j})_{1 \le i,j \le n}$ be an $n \times n$ matrix of indeterminates over $\mathbb{C}$, and let
\[
  \perm_n(X) = \sum_{\sigma \in S_n} \prod_{i=1}^n x_{i,\sigma(i)}, \qquad
  \det_n(X) = \sum_{\sigma \in S_n} \mathrm{sgn}(\sigma) \prod_{i=1}^n x_{i,\sigma(i)}
\]
denote the permanent and determinant polynomials. Following Valiant~\cite{Valiant1979}, the \emph{determinantal complexity} $\dc_{\mathbb{C}}(f)$ of a polynomial $f \in \mathbb{C}[x_1,\dots,x_N]$ is the smallest positive integer $m$ such that $f$ can be written as $f(X) = \det(A_0 + L(X))$, where $A_0 \in \mathbb{C}^{m\times m}$ is a constant matrix and $L(X) = \sum_{a=1}^N x_a A_a$ is a homogeneous linear matrix over $\mathbb{C}$.

For $n = 2$, $\dc_{\mathbb{C}}(\perm_2) = 2$ trivially. For $n = 3$, Grenet~\cite{Grenet2011} gave the general upper bound $\dc_{\mathbb{C}}(\perm_n) \le 2^n - 1$, which yields $\dc_{\mathbb{C}}(\perm_3) \le 7$, and Alper, Bogart, and Velasco~\cite{ABV2017} proved $\dc_{\mathbb{C}}(f) \ge \codim(\Sing(V(f))) + 1$ for any hypersurface of degree $d > 2$ with $\codim(\Sing(V(f))) \ge 5$, establishing $\dc_{\mathbb{C}}(\perm_3) = 7$.

For $n = 4$, Grenet's construction gives $\dc_{\mathbb{C}}(\perm_4) \le 15$, which Landsberg and Ressayre~\cite{LR2017} showed to be optimal among representations preserving a large subgroup of the symmetry group of $\perm_4$. Without symmetry assumptions, the quadratic Hessian bound of Mignon and Ressayre~\cite{MR2004} yields $\dc_{\mathbb{C}}(\perm_4) \ge n^2/2 = 8$, while the singularity codimension bound of~\cite{ABV2017} yields $\dc_{\mathbb{C}}(\perm_4) \ge 8 + 1 = 9$ once $\codim_{\mathbb{C}^{16}}(\Sing(V(\perm_4))) = 8$ is established. Over $\mathbb{R}$, Yabe~\cite{Yabe2015} proved $\dc_{\mathbb{R}}(\perm_n) \ge (n-1)^2 + 1$, giving $\dc_{\mathbb{R}}(\perm_4) \ge 10$ (a bound valid over $\mathbb{R}$ only).

\begin{remark}[Scope and Literature Priority]
In a November 2026 preprint, Sheshadri~\cite{Sheshadri2026} claims the exact value $\dc_{\mathbb{C}}(\perm_4) = 15$ (not yet peer-reviewed) via a normal-form and Hessian-inequality analysis, and in the abstract describes $10$ as the best known lower bound (without specifying the field). Accordingly, the present note does not claim $\dc_{\mathbb{C}}(\perm_4) \ge 10$ as a new lower bound; rather, it serves as a self-contained methodological exposition via origin-isotropy, Skew-Square vanishing, and second-order singularity codimensions.
\end{remark}

\begin{theorem}[Main Theorem]\label{thm:main}
Over the complex numbers,
\[
  \dc_{\mathbb{C}}(\perm_4) \ge 10.
\]
\end{theorem}

\paragraph{Structure of the proof and epistemic transparency.}
Apart from the computer certificate of Theorem~\ref{thm:sing1}, the algebraic steps of our proof are self-contained and proved by hand in Sections~\ref{sec:sing2}--\ref{sec:mainproof}. The proof relies on exactly \emph{one} computer-assisted certificate (Theorem~\ref{thm:sing1}): a direct degree-$9$ finite-field Macaulay rank computation on the explicit $8$-dimensional integer subspace $L_8 \subset \mathbb{C}^{16}$ (normatively specified in Appendix~\ref{app:ml8}) proving $\codim_{\mathbb{C}^{16}}(\Sing(V(\perm_4))) \ge 8$. The primary direct C rank engine (\texttt{exact\_mod\_rref.c}) uses pure \texttt{uint64\_t} modular integer arithmetic with zero floating-point operations; the auxiliary Python IMSC cross-check (\texttt{verify\_codim\_sing\_perm4.py}) uses \texttt{float64} BLAS matrix multiplication (\texttt{R\_canon @ C\_s}) for intermediate table compression, which remains exact as integer arithmetic only because all inner-product sums stay strictly below $2^{53}$. A standalone negative control on a $9$-dimensional subspace $L_9$, where the corank must remain strictly positive, is also documented in Section~\ref{sec:cert}.

\section{Preliminaries: Strength and Singularity Strata}\label{sec:sing2}

Let $F \in \mathbb{C}[x_1,\dots,x_N]_d$ be a homogeneous polynomial of degree $d \ge 2$.
For $k \ge 1$, the \emph{$k$-th order singular locus} of $V(F) \subset \mathbb{C}^N$ is the affine cone
\[
  \Sing_k(V(F)) := \left\{ X \in \mathbb{C}^N : \frac{\partial^{|\alpha|} F}{\partial X^\alpha}(X) = 0 \text{ for all multi-indices } \alpha \text{ with } |\alpha| = k \right\}.
\]
By Euler's homogeneous function theorem, $\Sing_k(V(F))$ is equivalently the locus of points where $F$ vanishes to order at least $k+1$. We write $\Sing(V(F)) := \Sing_1(V(F))$.

\begin{definition}
The \emph{strength} $\str(F)$ of a homogeneous polynomial $F \in \mathbb{C}[x_1,\dots,x_N]_d$ is the smallest integer $k \ge 1$ such that
\[
  F = \sum_{i=1}^k G_i H_i
\]
for homogeneous polynomials $G_i, H_i \in \mathbb{C}[x_1,\dots,x_N]$ with $1 \le \deg(G_i), \deg(H_i) \le d - 1$.
\end{definition}

\begin{lemma}[Strength--Singularity Inequality]\label{lem:strength}
For any non-zero homogeneous polynomial $F \in \mathbb{C}[x_1,\dots,x_N]_d$ ($d \ge 2$),
\[
  \codim_{\mathbb{C}^N}(\Sing(V(F))) \le 2\,\str(F).
\]
\end{lemma}

\begin{proof}
Write $F = \sum_{i=1}^k G_i H_i$ with $k = \str(F)$ and $\deg(G_i), \deg(H_i) \ge 1$. At any point $X \in Z := V(G_1,\dots,G_k, H_1,\dots,H_k)$, the product rule yields
\[
  \nabla F(X) = \sum_{i=1}^k \big(H_i(X)\nabla G_i(X) + G_i(X)\nabla H_i(X)\big) = 0,
\]
so $Z \subseteq \Sing(V(F))$. Since $0 \in Z$ and $Z$ is defined by $2k$ equations in $\mathbb{C}^N$, Krull's principal ideal theorem implies that every irreducible component of $Z$ has codimension at most $2k$. Hence $\codim_{\mathbb{C}^N}(\Sing(V(F))) \le \codim_{\mathbb{C}^N}(Z) \le 2k$.
\end{proof}

We next determine the exact codimension of $\Sing_2(V(\perm_4))$ completely by hand.

\begin{proposition}\label{prop:sing2}
Over $\mathbb{C}$ (in fact, over any field of characteristic $\neq 2$),
\[
  \codim_{\mathbb{C}^{16}}\big(\Sing_2(V(\perm_4))\big) = 12.
\]
\end{proposition}

\begin{proof}
Since $\perm_4(X)$ is multilinear in each row and column, its non-identically-zero second partial derivatives are precisely the $36$ $2 \times 2$ sub-permanents of $X$:
\[
  \frac{\partial^2 \perm_4}{\partial x_{i_1,j_1} \partial x_{i_2,j_2}}(X)
  = \perm_2\big(X_{\{1,2,3,4\}\setminus\{i_1,i_2\},\,\{1,2,3,4\}\setminus\{j_1,j_2\}}\big)
  \qquad (i_1 \neq i_2,\ j_1 \neq j_2).
\]
Thus $X \in \Sing_2(V(\perm_4))$ if and only if
\begin{equation}\label{eq:subperm2}
  x_{i,j} x_{k,\ell} + x_{i,\ell} x_{k,j} = 0 \qquad \text{for all } 1 \le i < k \le 4,\ 1 \le j < \ell \le 4.
\end{equation}
Let $G(X)$ be the bipartite graph on vertex sets $R = \{1,2,3,4\}$ and $C = \{1,2,3,4\}$ whose edge set is the support $\supp(X) = \{(i,j) : x_{i,j} \neq 0\}$.

\emph{Claim 1:} If $(i,j), (k,\ell) \in \supp(X)$ with $i \neq k$ and $j \neq \ell$, then both $(i,\ell)$ and $(k,j)$ belong to $\supp(X)$.  
Indeed, $x_{i,j} x_{k,\ell} \neq 0$ together with~\eqref{eq:subperm2} forces $x_{i,\ell} x_{k,j} = -x_{i,j} x_{k,\ell} \neq 0$.

By Claim~1, no two disjoint edges can lie in different connected components of $G(X)$, and every $2$-edge matching completes to a $4$-cycle. Consequently, $G(X)$ has at most one connected component containing edges, and that component is a complete bipartite graph $K_{a,b}$ supported on $I \times J$ for some $I \subseteq R$, $J \subseteq C$ with $|I| = a$, $|J| = b$.

\emph{Claim 2:} We cannot have both $a \ge 2$ and $b \ge 3$ (nor, by transpose, $a \ge 3$ and $b \ge 2$).  
Suppose for contradiction that $x_{i,j} \neq 0$ for all $i \in \{1,2\}$ and $j \in \{1,2,3\}$. Applying~\eqref{eq:subperm2} to the three pairs of columns $\{1,2\}$, $\{2,3\}$, $\{1,3\}$ gives
\[
  \frac{x_{2,1}}{x_{1,1}} = -\frac{x_{2,2}}{x_{1,2}}, \qquad
  \frac{x_{2,2}}{x_{1,2}} = -\frac{x_{2,3}}{x_{1,3}}, \qquad
  \frac{x_{2,1}}{x_{1,1}} = -\frac{x_{2,3}}{x_{1,3}}.
\]
Combining the first two yields $x_{2,1}/x_{1,1} = +x_{2,3}/x_{1,3}$, contradicting the third since $\mathrm{char}(\mathbb{C}) \neq 2$ and all entries are non-zero.

Therefore, every $X \in \Sing_2(V(\perm_4))$ belongs to the finite union of:
\begin{enumerate}
  \item $4$ single-row subspaces ($K_{1,4}$), each of dimension $4$;
  \item $4$ single-column subspaces ($K_{4,1}$), each of dimension $4$;
  \item $36$ subvarieties supported on a $2 \times 2$ block $I \times J$ satisfying $\perm_2(X_{I,J}) = 0$, each of dimension $4 - 1 = 3$.
\end{enumerate}
Conversely, every matrix in any of these three families trivially satisfies~\eqref{eq:subperm2}. Hence $\dim \Sing_2(V(\perm_4)) = 4$, and $\codim_{\mathbb{C}^{16}}(\Sing_2(V(\perm_4))) = 16 - 4 = 12$.
\end{proof}

\section{First-Order Singularity Codimension via Macaulay Certificate}\label{sec:cert}

\begin{theorem}[Computer-Assisted Exact Certificate]\label{thm:sing1}
Over $\mathbb{C}$,
\[
  \codim_{\mathbb{C}^{16}}\big(\Sing(V(\perm_4))\big) = 8,
\]
and consequently $\str(\perm_4) = 4$.
\end{theorem}

\begin{proof}
For the upper bound $\codim_{\mathbb{C}^{16}}(\Sing(V(\perm_4))) \le 8$, observe that any $4 \times 4$ matrix whose first two rows are zero has all $3 \times 3$ sub-permanents equal to $0$. This gives an $8$-dimensional linear subspace contained in $\Sing(V(\perm_4))$, so $\dim \Sing(V(\perm_4)) \ge 8$.

For the lower bound $\codim_{\mathbb{C}^{16}}(\Sing(V(\perm_4))) \ge 8$, let $J = (\partial_{i,j}\perm_4)_{1\le i,j\le 4} \subset \mathbb{Z}[X]$ be the Jacobian ideal of $\perm_4$, generated by the $16$ cubic $3 \times 3$ sub-permanents. Fix an explicit $16 \times 8$ integer matrix $M_{L_8}$ of rank $8$, defining an $8$-dimensional linear subspace $L_8 = \{M_{L_8} t : t \in \mathbb{C}^8\} \subset \mathbb{C}^{16}$. Let $I := J|_{L_8} \subset R := \mathbb{Z}[t_1,\dots,t_8]$ be the restricted ideal.

To prove $V_{\mathbb{C}}(I) = \{0\}$, it suffices by the homogeneous Nullstellensatz to exhibit a degree $D$ such that $I_D \otimes \mathbb{C} = R_D \otimes \mathbb{C}$, i.e., such that the degree-$D$ Macaulay matrix $M_D(I)$ (whose rows are the coefficient vectors of $m \cdot g$ for all monomials $m \in R_{D-3}$ and all $16$ generators $g$ of $I$) has full column rank $\binom{D+7}{7}$ over $\mathbb{Q}$. Since $M_D(I)$ has integer entries, if $M_D(I) \bmod p$ has full column rank $\binom{D+7}{7}$ over $\mathbb{F}_p$ for some prime $p$, then $M_D(I)$ has a non-vanishing maximal minor modulo $p$, hence non-vanishing over $\mathbb{Z}$ and $\mathbb{C}$.

\emph{Primary Direct Certificate ($D = 9$, no IMSC compression):} For the explicit $16 \times 8$ integer matrix $M_{L_8}$ given in Appendix~\ref{app:ml8} (generated by \texttt{np.random.default\_rng(2026).integers(-2, 3, size=(16, 8))}, rank~$8$), we expand the $16$ restricted cubic generators in $\mathbb{Z}[t_1,\dots,t_8]$ (independently cross-checked against direct $S_3$ permutation-sum evaluation at $5$ random points in $\mathbb{F}_{32{,}749}^8$, passing $80/80$ checks) and form the full degree-$9$ Macaulay matrix $M_9(I)$ directly in CSR format ($27{,}456$ rows $= 16 \times \binom{13}{6}$, $11{,}440$ columns $= \binom{16}{7}$, $3{,}003{,}000$ nonzeros). Exact modular row reduction via the standalone pure-\texttt{uint64\_t} C engine \texttt{exact\_mod\_rref.c} modulo $p_1 = 32{,}749$ yields full column rank:
\[
  \rank_{\mathbb{F}_{32{,}749}} M_9(I) = 11{,}440 / 11{,}440 \qquad (\text{corank} = 0).
\]
Note that this primary direct $D = 9$ certificate is evaluated modulo a single prime ($p_1 = 32{,}749$, which suffices to certify non-vanishing of a maximal minor over $\mathbb{Z}$ and $\mathbb{C}$), whereas the IMSC cross-check below is evaluated modulo two primes ($p_1 = 32{,}749$ and $p_2 = 65{,}521$).

\emph{Independent IMSC Cross-Check:} As an independent cross-check, because $I$ is generated in degree $3$, for every $D \ge 6$ we have $I_{D+1} = \sum_{i=1}^8 t_i I_D$. Let $\mathcal{B}_D$ be a monomial basis of $(R_D / I_D) \otimes \mathbb{F}_p$, with reduction table $C_D : R_D \to \Span_{\mathbb{F}_p}(\mathcal{B}_D)$ satisfying $m \equiv C_D(m) \pmod{I_D}$ for all $m \in R_D$. Then $S_{D+1} := \{t_i b : 1 \le i \le 8,\ b \in \mathcal{B}_D\}$ spans $(R_{D+1}/I_{D+1}) \otimes \mathbb{F}_p$. For every monomial $\beta \in R_{D+1}$ and every pair of factorizations $\beta = t_i \alpha_i = t_j \alpha_j$, the difference $t_i C_D(\alpha_i) - t_j C_D(\alpha_j)$ is a linear combination of elements of $S_{D+1}$ that lies in $I_{D+1} \otimes \mathbb{F}_p$. Executing this Iterated Macaulay Schur-Complement (IMSC) cascade in \texttt{verify\_codim\_sing\_perm4.py} on the same $L_8$ yields, identically modulo both $p_1 = 32{,}749$ and $p_2 = 65{,}521$:
\[
\begin{array}{cccc}
\toprule
\text{Degree } D & \dim R_D = \binom{D+7}{7} & \dim(I_D \otimes \mathbb{F}_p) & \text{corank} \\
\midrule
6 & 1{,}716 & 1{,}584 & 132 \\
7 & 3{,}432 & 3{,}400 & 32 \\
8 & 6{,}435 & 6{,}427 & 8 \\
9 & 11{,}440 & 11{,}440 & 0 \\
\bottomrule
\end{array}
\]
At $D = 9$, the corank vanishes identically, so $\Sing(V(\perm_4)) \cap L_8 = \{0\}$. Since $\Sing(V(\perm_4))$ is a homogeneous cone, $0 \in Y_i \cap L_8$ for every irreducible component $Y_i$ of $\Sing(V(\perm_4))$; by the affine dimension theorem, $\dim Y_i + 8 - 16 \le \dim(Y_i \cap L_8) = 0$ for every component $Y_i$, giving $\codim_{\mathbb{C}^{16}}(\Sing(V(\perm_4))) \ge 8$, and hence $= 8$.

Finally, Lemma~\ref{lem:strength} gives $\str(\perm_4) \ge \lceil 8/2 \rceil = 4$, while Laplace expansion along the first row, $\perm_4(X) = \sum_{j=1}^4 x_{1,j}\,\perm_3(X_{\hat{1},\hat{j}})$, gives $\str(\perm_4) \le 4$. Thus $\str(\perm_4) = 4$.
\end{proof}

\begin{remark}[Negative Control on $L_9$]
Since $\dim \Sing(V(\perm_4)) = 8$, a generic $9$-dimensional subspace $L_9 \subset \mathbb{C}^{16}$ must intersect $\Sing(V(\perm_4))$ in a cone of dimension at least $1$, so the corank of $J|_{L_9}$ must remain strictly positive in every degree. As a sanity check against implementation errors in the IMSC routine, \texttt{verify\_codim\_sing\_perm4.py} runs the identical algorithm on an integer $9$-dimensional subspace $L_9$ modulo $p_1 = 32{,}749$, obtaining strictly positive coranks $753$ at $D = 6$ and $580$ at $D = 7$.
\end{remark}

\section{Origin-Corank Classification}\label{sec:corank}

Let $\perm_4(X) = \det(M(X))$ be an affine determinantal representation of size $m$, where $M(X) = A_0 + L(X)$ with $A_0 \in \mathbb{C}^{m\times m}$ constant and $L(X)$ homogeneous linear in $X \in \mathbb{C}^{16}$. Let $r := \corank(A_0) = m - \rank(A_0)$ and $w := \rank(A_0) = m - r$. Multiplying $M(X)$ on the left and right by constant invertible matrices $P, Q \in \mathrm{GL}_m(\mathbb{C})$ with $\det(PQ) = 1$, we may assume
\begin{equation}\label{eq:blockM}
  M(tX) = A_0 + tL(X) = \begin{pmatrix} t B(X) & t U(X) \\ t V(X) & I_w + t A(X) \end{pmatrix},
\end{equation}
where $B \in \mathbb{C}[X]_1^{r\times r}$, $U \in \mathbb{C}[X]_1^{r\times w}$, $V \in \mathbb{C}[X]_1^{w\times r}$, and $A \in \mathbb{C}[X]_1^{w\times w}$ are homogeneous linear matrices. For $t$ in a neighborhood of $0$, $I_w + tA(X)$ is invertible, and the Schur complement identity yields
\begin{equation}\label{eq:schur}
  t^4 \perm_4(X) = \det(M(tX))
  = t^r \det(I_w + tA) \cdot \det\!\left(B - \sum_{k=0}^{\infty} (-1)^k t^{k+1} U A^k V\right).
\end{equation}

\begin{proposition}[Mandatory Corank $1$]\label{prop:corank1}
Every affine determinantal representation $\perm_4(X) = \det(A_0 + L(X))$ over $\mathbb{C}$ satisfies $\corank(A_0) = 1$.
\end{proposition}

\begin{proof}
We eliminate all $r \neq 1$ case by case using Theorem~\ref{thm:sing1} ($\codim(\Sing(V(\perm_4))) = 8$ and $\str(\perm_4) = 4$).

\emph{Case $r = 0$.} Evaluating $M(X) = A_0 + L(X)$ at $X = 0$ gives $0 = \perm_4(0) = \det(A_0)$, contradicting $\rank(A_0) = m$ ($r = 0$).

\emph{Case $r \ge 5$.} By~\eqref{eq:schur}, $\det(M(tX))$ is divisible by $t^r$ with $r \ge 5$, so $[t^4]\det(M(tX)) = 0 \neq \perm_4(X)$, impossible.

\emph{Case $r = 4$.} Comparing $[t^4]$ in~\eqref{eq:schur} gives $\perm_4(X) = \det_4(B(X))$ for a $4\times 4$ linear matrix $B(X)$. Let $\Sigma_2 \subset \mathbb{C}^{4\times 4}$ be the variety of $4\times 4$ matrices of rank $\le 2$, which has codimension $(4-2)^2 = 4$. At every $X \in B^{-1}(\Sigma_2)$, all $3\times 3$ minors of $B(X)$ vanish, so $\nabla(\det_4(B(X))) = 0$, i.e., $B^{-1}(\Sigma_2) \subseteq \Sing(V(\perm_4))$. Since $0 \in B^{-1}(\Sigma_2)$, we obtain $\codim_{\mathbb{C}^{16}}(\Sing(V(\perm_4))) \le \codim_{\mathbb{C}^{16}}(B^{-1}(\Sigma_2)) \le 4 < 8$, contradicting Theorem~\ref{thm:sing1}.

\emph{Case $r = 3$.} Comparing $[t^3]$ and $[t^4]$ in~\eqref{eq:schur} gives $\det_3(B(X)) \equiv 0$ and
\[
  \perm_4(X) = \tr(A)\underbrace{\det_3(B)}_{=\,0} - \tr\!\big(\adj(B(X))\, U(X) V(X)\big)
  = -\tr\!\big(\adj(B(X))\, U(X) V(X)\big).
\]
Thus $B(X)$ spans a linear subspace of $\mathbb{C}^{3\times 3}$ consisting of matrices of rank $\le 2$. Left and right multiplication of $M(tX)$ in~\eqref{eq:blockM} by $\mathrm{diag}(P_1, I_w)$ and $\mathrm{diag}(Q_1, I_w)$ ($P_1, Q_1 \in \mathrm{GL}_3(\mathbb{C})$) replaces $(B, U, V)$ by $(P_1 B Q_1, P_1 U, V Q_1)$ while scaling $\det M(tX)$ by $c := \det(P_1 Q_1) \in \mathbb{C}^\times$; transposition $M(tX) \mapsto M(tX)^T$ preserves $\det M(tX)$ and replaces $(B, U, V, A)$ by $(B^T, V^T, U^T, A^T)$. Denoting the resulting blocks again by $(B, U, V, A)$, we still have $c\,\perm_4(X) = -\tr(\adj(B(X)) U(X) V(X))$. Up to these operations, every linear space of $3 \times 3$ matrices of rank $\le 2$ over $\mathbb{C}$ falls into one of three types (cf.\ Atkinson~\cite{Atkinson1983}, Eisenbud--Harris~\cite{EH1988}; a self-contained verification is given immediately below):
\begin{enumerate}
  \item \emph{Compression type I ($1\times 3$ zero row, or $3\times 1$ zero column via transpose):} $B(X)$ has a zero third row. Then $\adj(B(X)) = c(X) e_3^T$, where $c(X) = (c_1,c_2,c_3)^T \in \mathbb{C}[X]_2^3$ consists of the signed $2\times 2$ minors of the top $2\times 3$ block of $B(X)$. Hence
  \[
    c\,\perm_4(X) = -\big(e_3^T U(X) V(X)\big) c(X) = -\sum_{j=1}^3 \big(e_3^T U(X) V(X)\big)_j\, c_j(X),
  \]
  giving $\str(\perm_4) = \str(c\,\perm_4) \le 3 < 4$, a contradiction.
  \item \emph{Compression type II ($2\times 2$ zero block):} Permuting rows and columns puts the zero block in rows $\{2,3\}$ and columns $\{2,3\}$, so
  \[
    B(X) = \begin{pmatrix} b_{11}(X) & b_{12}(X) & b_{13}(X) \\ b_{21}(X) & 0 & 0 \\ b_{31}(X) & 0 & 0 \end{pmatrix}.
  \]
  All $2\times 2$ minors involving two rows from $\{2,3\}$ or two columns from $\{2,3\}$ vanish, so the first row and first column of $\adj(B(X))$ are zero, and computing the remaining four cofactors gives the rank-$\le 1$ outer product
  \[
    \adj(B(X)) = p(X)\, q(X)^T, \qquad
    p(X) := \begin{pmatrix} 0 \\ -b_{13}(X) \\ b_{12}(X) \end{pmatrix}, \quad
    q(X) := \begin{pmatrix} 0 \\ b_{31}(X) \\ -b_{21}(X) \end{pmatrix}.
  \]
  Consequently,
  \[
    c\,\perm_4(X) = -q(X)^T U(X) V(X) p(X)
    = -b_{31}(X)\big(U(X)V(X)p(X)\big)_2 + b_{21}(X)\big(U(X)V(X)p(X)\big)_3,
  \]
  which gives $\str(\perm_4) = \str(c\,\perm_4) \le 2 < 4$, a contradiction.
  \item \emph{Skew-symmetric primitive type:} $B(X) \in \mathfrak{so}_3(\mathbb{C})$, i.e.,
  \[
    B(X) = \begin{pmatrix} 0 & \ell_3 & -\ell_2 \\ -\ell_3 & 0 & \ell_1 \\ \ell_2 & -\ell_1 & 0 \end{pmatrix}
  \]
  for linear forms $\ell_1,\ell_2,\ell_3 \in \mathbb{C}[X]_1$. Then $\adj(B(X)) = \ell(X)\ell(X)^T$, so
  \[
    c\,\perm_4(X) = -\ell(X)^T U(X) V(X) \ell(X) = -\sum_{i=1}^3 \ell_i(X)\, \big(U(X) V(X) \ell(X)\big)_i,
  \]
  again giving $\str(\perm_4) = \str(c\,\perm_4) \le 3 < 4$, a contradiction.
\end{enumerate}

\emph{Self-contained verification of the $3\times 3$ classification:} Since $c\,\perm_4 \neq 0$, we have $B \not\equiv 0$. If $\max_{X \in \mathbb{C}^{16}} \rank B(X) = 1$, choose $X_0$ with $B(X_0) = E_{11}$ up to $\mathrm{GL}_3(\mathbb{C})\times\mathrm{GL}_3(\mathbb{C})$; vanishing of all $2\times 2$ minors of $t E_{11} + B(X)$ forces $b_{ij} \equiv 0$ ($i,j\in\{2,3\}$) and $b_{1j} b_{i1} \equiv 0$, so in the integral domain $\mathbb{C}[X]$ either rows $\{2,3\}$ or columns $\{2,3\}$ vanish (Type~I). Otherwise $\max_X \rank B(X) = 2$; choose $X_0$ with $B(X_0) = \mathrm{diag}(1,1,0)$ and write $B(X) = \left(\begin{smallmatrix} P(X) & q(X) \\ r(X)^T & s(X) \end{smallmatrix}\right)$ with $P \in \mathbb{C}[X]_1^{2\times 2}$ and $q,r \in \mathbb{C}[X]_1^2$. From $\det_3(t B(X_0) + B(X)) \equiv 0$, $[t^2] = 0$ gives $s \equiv 0$ and $[t^1] = 0$ gives $r^T q \equiv 0$, which polarizes to $r(X)^T q(Y) + r(Y)^T q(X) \equiv 0$. Let $Q := \Span_{\mathbb{C}}\{q(X)\} \subseteq \mathbb{C}^2$ and $R := \Span_{\mathbb{C}}\{r(X)\} \subseteq \mathbb{C}^2$. If $\dim R = 0$ or $\dim Q = 0$, $B$ is of Type~I. If $\dim Q = \dim R = 1$, conjugating by $\mathrm{diag}(g,1)$ and $\mathrm{diag}(g^{-1},1)$ ($g \in \mathrm{GL}_2(\mathbb{C})$) preserves $B(X_0)$ and sets $q_2 \equiv 0$ with $q_1 \not\equiv 0$; polarity forces $r_1 \equiv 0$ with $r_2 \not\equiv 0$, and then $0 \equiv \det_3(B(X)) = q_1 r_2 p_{21}$ forces $p_{21} \equiv 0$, placing $B$ in Type~II (rows $\{2,3\}$ and columns $\{1,3\}$ zero). Finally, if $\dim Q = 2$ (or $\dim R = 2$, reduced to $\dim Q = 2$ by transposition), polarity $r(X)^T q(Y) + r(Y)^T q(X) \equiv 0$ implies $q(X) = 0 \implies r(X) = 0$ (since $Q = \mathbb{C}^2$). Hence $r(X) = K_2 q(X)$ for a $2\times 2$ matrix $K_2$ satisfying $K_2^T = -K_2$. Since $K_2 = 0$ gives $r \equiv 0$ (Type~I), we have $K_2 \neq 0$; rescaling row and column~$3$ gives $r = (-q_2, q_1)^T$, and cofactor expansion along row~$3$ yields
\[
  0 \equiv \det_3(B(X)) = p_{21} q_1^2 + (p_{22} - p_{11}) q_1 q_2 - p_{12} q_2^2.
\]
Decomposing $\mathbb{C}^{16} = \ker(q) \oplus V_2$ with coordinates $(Y, z_1, z_2)$ where $Y \in \ker(q)$ and $z_i = q_i(X)$, we write $p_{ij}(X) = p_{ij}(Y) + \alpha_{ij} z_1 + \beta_{ij} z_2$. Collecting terms linear in $Y$ in $\det_3(B(X)) \equiv 0$ forces $p_{21}(Y) \equiv p_{12}(Y) \equiv p_{22}(Y) - p_{11}(Y) \equiv 0$ on $\ker(q)$, while the cubic terms in $(z_1, z_2)$ give $\alpha_{21} = \beta_{12} = 0$, $\alpha_{22} - \alpha_{11} = -\beta_{21}$, and $\beta_{22} - \beta_{11} = \alpha_{12}$. Setting $\alpha := \beta_{21}$ and $\gamma := \alpha_{12}$, we obtain $p_{21} = \alpha q_2$, $p_{12} = \gamma q_1$, and $p_{22} - p_{11} = -\alpha q_1 + \gamma q_2$. Direct multiplication then gives
\[
  \begin{pmatrix} 1 & 0 & \alpha \\ 0 & 1 & \gamma \\ 0 & 0 & 1 \end{pmatrix}
  \begin{pmatrix} 0 & -1 & 0 \\ 1 & 0 & 0 \\ 0 & 0 & -1 \end{pmatrix}
  B(X)
  =
  \begin{pmatrix} 1 & 0 & \alpha \\ 0 & 1 & \gamma \\ 0 & 0 & 1 \end{pmatrix}
  \begin{pmatrix} -p_{21} & -p_{22} & -q_2 \\ p_{11} & p_{12} & q_1 \\ q_2 & -q_1 & 0 \end{pmatrix}
  =
  \begin{pmatrix} 0 & -(p_{11} + \gamma q_2) & -q_2 \\ p_{11} + \gamma q_2 & 0 & q_1 \\ q_2 & -q_1 & 0 \end{pmatrix}
  \in \mathfrak{so}_3(\mathbb{C})
\]
(using $p_{21} = \alpha q_2$, $p_{12} = \gamma q_1$, and $p_{22} - p_{11} = -\alpha q_1 + \gamma q_2$), placing $B(X)$ in Type~III.

\emph{Case $r = 2$.} Let $C_1 = -U(X)V(X)$ and $C_2 = U(X)A(X)V(X)$ (both $2\times 2$). Expanding $\det(B + t C_1 + t^2 C_2) = d_2 + t d_3 + t^2 d_4 + \cdots$ in~\eqref{eq:schur} (where the cross terms $\tr(A)d_2$ in $[t^3]$ and $\tr(A)d_3, e_2(A)d_2$ in $[t^4]$ vanish sequentially once $d_2 \equiv 0$ and $d_3 \equiv 0$) yields:
\begin{align*}
  [t^2] = 0 &\implies d_2 = \det_2(B(X)) \equiv 0, \\
  [t^3] = 0 &\implies d_3 = -\tr\!\big(\adj(B(X))\, U(X) V(X)\big) \equiv 0, \\
  [t^4] = \perm_4(X) &\implies \perm_4(X) = d_4 = \det_2(U(X)V(X)) + \tr\!\big(\adj(B(X))\, U(X) A(X) V(X)\big).
\end{align*}
Since $\det_2(B(X)) \equiv 0$, $B(X)$ is a $2\times 2$ linear matrix of rank $\le 1$. Up to left/right $\mathrm{GL}_2(\mathbb{C})$ action and transposition, $B(X)$ has a zero second row:
\[
  B(X) = \begin{pmatrix} b_1(X) & b_2(X) \\ 0 & 0 \end{pmatrix}, \qquad
  \adj(B(X)) = \begin{pmatrix} 0 & -b_2(X) \\ 0 & b_1(X) \end{pmatrix}.
\]
Let $Q = U(X)V(X)$ and $R = U(X)A(X)V(X)$. The $[t^3] = 0$ condition reads
\[
  b_1(X)\, Q_{22}(X) = b_2(X)\, Q_{21}(X),
\]
and the $[t^4]$ identity reads
\[
  \perm_4(X) = Q_{11} Q_{22} - Q_{12} Q_{21} - b_2 R_{21} + b_1 R_{22}.
\]
\begin{itemize}
  \item If $b_1 = b_2 = 0$, then $\perm_4 = Q_{11} Q_{22} - Q_{12} Q_{21}$, so $\str(\perm_4) \le 2 < 4$.
  \item If $b_1, b_2$ are linearly dependent and not both zero, a column operation sets $b_2 = 0$, $b_1 \neq 0$. Then $b_1 Q_{22} = 0$ forces $Q_{22} \equiv 0$ (since $\mathbb{C}[X]$ is an integral domain), leaving $\perm_4 = -Q_{12} Q_{21} + b_1 R_{22}$, so $\str(\perm_4) \le 2 < 4$.
  \item If $b_1, b_2$ are linearly independent, then $\gcd(b_1, b_2) = 1$ in the UFD $\mathbb{C}[X]$, so $b_1 Q_{22} = b_2 Q_{21}$ forces $Q_{21} = b_1 \ell$ and $Q_{22} = b_2 \ell$ for some linear form $\ell \in \mathbb{C}[X]_1$. Substituting into $\perm_4$ yields
  \[
    \perm_4(X) = b_1\big(R_{22} - \ell Q_{12}\big) + b_2\big(\ell Q_{11} - R_{21}\big),
  \]
  again giving $\str(\perm_4) \le 2 < 4$, a contradiction.
\end{itemize}
Hence $r = 1$ is the only remaining possibility.
\end{proof}

\section{Origin-Isotropy, Skew-Square Vanishing, and Proof of Theorem~\ref{thm:main}}\label{sec:mainproof}

We first record the elementary matrix identity at the heart of our proof.

\begin{lemma}[Skew-Square Lemma]\label{lem:skewsquare}
Let $M, K \in \mathbb{C}^{d \times d}$ satisfy $K^T = -K$ and $(MK)^T = -MK$. Then
\[
  (M^2 K)^T = -M^2 K.
\]
In particular, $y^T M^2 K y = 0$ for all $y \in \mathbb{C}^d$.
\end{lemma}

\begin{proof}
Since $K^T = -K$, the hypothesis $(MK)^T = -MK$ reads $-K M^T = K^T M^T = -MK$, hence $K M^T = MK$. Therefore,
\[
  (M^2 K)^T = K^T (M^T)^2 = (-K M^T) M^T = (-MK) M^T = -M(K M^T) = -M(MK) = -M^2 K.
\]
Note that $K$ is not assumed to be invertible. Because $M^2 K$ is skew-symmetric, $y^T M^2 K y = 0$ identically.
\end{proof}

\begin{proof}[Proof of Theorem~\ref{thm:main}]
Suppose for contradiction that $\dc_{\mathbb{C}}(\perm_4) \le 9$. By padding with an identity block if necessary, there exists a $9 \times 9$ affine determinantal representation $\perm_4(X) = \det(A_0 + L(X))$. By Proposition~\ref{prop:corank1}, $\corank(A_0) = 1$, so $w = \rank(A_0) = 8$. Writing $v(X)^T = U(X) \in \mathbb{C}[X]_1^{1\times 8}$ and $u(X) = V(X) \in \mathbb{C}[X]_1^{8\times 1}$ in~\eqref{eq:blockM}--\eqref{eq:schur}, we have
\[
  t^4 \perm_4(X) = t \det(I_8 + tA(X))\Big(b(X) - t\,v(X)^T u(X) + t^2\,v(X)^T A(X) u(X) - t^3\,v(X)^T A(X)^2 u(X) + \cdots\Big).
\]
Comparing powers of $t$ sequentially yields:
\begin{align}
  [t^1] = 0 &\implies b(X) \equiv 0, \nonumber\\
  [t^2] = 0 &\implies v(X)^T u(X) \equiv 0, \label{eq:t2}\\
  [t^3] = 0 &\implies v(X)^T A(X) u(X) \equiv 0, \label{eq:t3}\\
  [t^4] = \perm_4(X) &\implies \perm_4(X) = -v(X)^T A(X)^2 u(X). \label{eq:t4}
\end{align}
Replacing $v$ by $-v$, we may drop the minus sign in~\eqref{eq:t4} without affecting~\eqref{eq:t2}--\eqref{eq:t3}.

\paragraph{Step 1: Origin-isotropy forces $\dim W = 8$.}
Let $W := \Span_{\mathbb{C}}\{v_1,\dots,v_8, u_1,\dots,u_8\} \subseteq (\mathbb{C}^{16})^*$ and $D := \dim W$. Consider the non-degenerate symmetric bilinear form $\Omega$ on $\mathbb{C}^{16} = \mathbb{C}^8 \oplus \mathbb{C}^8$ defined by
\[
  \Omega\!\left(\begin{pmatrix} \mathbf{v} \\ \mathbf{u} \end{pmatrix}, \begin{pmatrix} \mathbf{v}' \\ \mathbf{u}' \end{pmatrix}\right) = \mathbf{v}^T \mathbf{u}' + (\mathbf{v}')^T \mathbf{u}.
\]
Polarizing~\eqref{eq:t2} gives $v(X)^T u(Y) + v(Y)^T u(X) = 0$ for all $X, Y \in \mathbb{C}^{16}$. Hence the image subspace $S := \{(v(X), u(X))^T : X \in \mathbb{C}^{16}\} \subseteq \mathbb{C}^{16}$ is totally isotropic with respect to $\Omega$, which implies
\[
  D = \dim W = \dim S \le 8.
\]
On the other hand, on $W^\perp := \{X \in \mathbb{C}^{16} : v(X) = 0 \text{ and } u(X) = 0\}$,~\eqref{eq:t4} shows that $\perm_4$ belongs to the square of the defining ideal $I_W$ of $W^\perp$, so $W^\perp \subseteq \Sing(V(\perm_4))$. By Theorem~\ref{thm:sing1},
\[
  8 = \codim_{\mathbb{C}^{16}}(\Sing(V(\perm_4))) \le \codim_{\mathbb{C}^{16}}(W^\perp) = D \le 8,
\]
and therefore $D = 8$.

\paragraph{Step 2: $2$-block normal form.}
For any $P \in \mathrm{GL}_8(\mathbb{C})$, the simultaneous conjugation $(v^T, A, u) \mapsto (v^T P^{-1}, P A P^{-1}, P u)$ preserves~\eqref{eq:t2}--\eqref{eq:t4}. Let $d_v := \dim \Span\{v_1,\dots,v_8\} \ge 1$ (since $v \equiv 0$ would contradict~\eqref{eq:t4}) and $d_z := 8 - d_v \ge 0$ (treating the second block as empty if $d_z = 0$). Choosing $P \in \mathrm{GL}_8(\mathbb{C})$ so that the first $d_v$ entries of $v^T P^{-1}$ form a basis $y = (y_1,\dots,y_{d_v})^T$ of $\Span\{v_1,\dots,v_8\}$ and the remaining $d_z$ entries vanish, we obtain
\[
  v(X) = \begin{pmatrix} y(X) \\ 0 \end{pmatrix}, \qquad
  u(X) = \begin{pmatrix} u_1(X) \\ z(X) \end{pmatrix}, \qquad
  A(X) = \begin{pmatrix} A_{11}(X) & A_{12}(X) \\ A_{21}(X) & A_{22}(X) \end{pmatrix},
\]
with $u_1 \in \mathbb{C}[X]_1^{d_v}$ and $z \in \mathbb{C}[X]_1^{d_z}$.

Condition~\eqref{eq:t2} becomes $y(X)^T u_1(X) \equiv 0$. Extending $y_1,\dots,y_{d_v}$ to a basis of $(\mathbb{C}^{16})^*$ and writing $u_1(X) = K y(X) + \tilde{u}_1$ (where $\tilde{u}_1$ depends only on the complementary basis elements), comparing coefficients of mixed products forces $\tilde{u}_1 \equiv 0$ and comparing quadratic terms in $y$ forces $K^T = -K$. Thus $u_1(X) = K y(X)$ with $K \in \mathfrak{so}_{d_v}(\mathbb{C})$.

Because $u_1 = K y$ already lies in $\Span(y_1,\dots,y_{d_v})$, we have $W = \Span\{y_1,\dots,y_{d_v}, z_1,\dots,z_{d_z}\}$. Since $\dim W = 8 = d_v + d_z$, the $8$ linear forms $(y_1,\dots,y_{d_v}, z_1,\dots,z_{d_z})$ are linearly independent and form a basis of $W$.

\paragraph{Step 3: Extracting the $\xi$-linear constraints from $[t^3] = 0$.}
Extend $(y_1,\dots,y_{d_v}, z_1,\dots,z_{d_z})$ to a basis $(y_1,\dots,y_{d_v}, z_1,\dots,z_{d_z}, \xi_1,\dots,\xi_8)$ of $(\mathbb{C}^{16})^*$. Then $I_W = \langle y_1,\dots,y_{d_v}, z_1,\dots,z_{d_z} \rangle$, and $W^\perp = V(I_W)$ is parametrized by the $8$ free coordinates $\xi = (\xi_1,\dots,\xi_8)$. Decompose
\[
  A_{11}(X) = A_{11}^{(\xi)}(\xi) + A_{11}^{(W)}(y,z), \qquad
  A_{12}(X) = A_{12}^{(\xi)}(\xi) + A_{12}^{(W)}(y,z),
\]
where the entries of $A_{11}^{(W)}, A_{12}^{(W)}$ lie in $I_W$. Substituting into~\eqref{eq:t3} gives
\[
  0 \equiv v(X)^T A(X) u(X) = y^T A_{11}(X) K y + y^T A_{12}(X) z.
\]
Collecting the terms of degree $1$ in $\xi$ (which must vanish identically as a polynomial in the independent variables $y, z, \xi$), we obtain
\[
  y^T A_{11}^{(\xi)}(\xi) K y + y^T A_{12}^{(\xi)}(\xi) z \equiv 0.
\]
Separating by bidegree in $(y, z)$:
\begin{enumerate}
  \item The $(1,1)$-bidegree part $y^T A_{12}^{(\xi)}(\xi) z \equiv 0$ forces $A_{12}^{(\xi)}(\xi) \equiv 0$, i.e., every entry of $A_{12}(X)$ belongs to $I_W$.
  \item The $(2,0)$-bidegree part $y^T A_{11}^{(\xi)}(\xi) K y \equiv 0$ forces the matrix $A_{11}^{(\xi)}(\xi) K$ to be skew-symmetric for every $\xi \in \mathbb{C}^8$:
  \[
    \big(A_{11}^{(\xi)}(\xi) K\big)^T = -A_{11}^{(\xi)}(\xi) K.
  \]
\end{enumerate}

\paragraph{Step 4: Order-$3$ vanishing along $W^\perp$ and contradiction.}
By Lemma~\ref{lem:skewsquare} applied to $M = A_{11}^{(\xi)}(\xi)$ and $K$, the matrix $A_{11}^{(\xi)}(\xi)^2 K$ is skew-symmetric for all $\xi \in \mathbb{C}^8$, so
\[
  y^T A_{11}^{(\xi)}(\xi)^2 K y \equiv 0.
\]
Now expand~\eqref{eq:t4} in our $2$-block partition:
\[
  \perm_4(X) = y^T A_{11}(X)^2 K y + y^T A_{12}(X) A_{21}(X) K y + y^T A_{11}(X) A_{12}(X) z + y^T A_{12}(X) A_{22}(X) z.
\]
Since $y, z \in I_W$ and $A_{12}(X) \in I_W$, the last three terms manifestly lie in $I_W^3$. For the first term,
\[
  A_{11}(X)^2 = A_{11}^{(\xi)}(\xi)^2 + \Big(A_{11}^{(\xi)}(\xi) A_{11}^{(W)} + A_{11}^{(W)} A_{11}^{(\xi)}(\xi) + \big(A_{11}^{(W)}\big)^2\Big),
\]
where the bracketed part lies in $I_W$. Since $y^T A_{11}^{(\xi)}(\xi)^2 K y \equiv 0$, it follows that $y^T A_{11}(X)^2 K y \in I_W^3$ as well. Hence $\perm_4(X) \in I_W^3$.

Consequently, all first and second partial derivatives of $\perm_4$ vanish identically on $W^\perp = V(I_W)$, so $W^\perp \subseteq \Sing_2(V(\perm_4))$. By Proposition~\ref{prop:sing2},
\[
  12 = \codim_{\mathbb{C}^{16}}(\Sing_2(V(\perm_4))) \le \codim_{\mathbb{C}^{16}}(W^\perp) = 8,
\]
a contradiction. Therefore $\dc_{\mathbb{C}}(\perm_4) \ge 10$.
\end{proof}

\begin{remark}[Scope and limitations]
Extending this normal-form approach to $m \ge 10$ would require introducing a middle block of size $w_{\mathrm{mid}} = w - D$ (bounding $\rank(\mathrm{Hess}(\perm_4)(\xi)) \le 2 w_{\mathrm{mid}}$ on $W^\perp$) and controlling the higher Hessian rank strata $\mathcal{H}_{\le 2k}(\perm_4)$---a possible direction that is not pursued here. As with all singularity- and Hessian-based techniques, this framework is subject to the quadratic Mignon--Ressayre barrier $O(n^2)$ as $n \to \infty$ and does not by itself address super-polynomial lower bounds for $\mathrm{VP} \neq \mathrm{VNP}$.
\end{remark}

\section{General Theory: Free-Algebra Word-Reversal, ABP/Hankel Obstructions, and the Degree-$n \ge 4$ Hypersurface Theorem}\label{sec:general}

In this section we generalize Lemma~\ref{lem:skewsquare} and the proof of Theorem~\ref{thm:main} in three stages: (1) a necessary and sufficient Word-Reversal and Darboux block characterization over the free associative algebra $\mathcal{A} = \mathbb{C}\langle \mathbf{m}_1,\dots,\mathbf{m}_s \rangle$; (2) width and rank obstructions for Algebraic Branching Programs (ABPs) and LTI Hankel matrices; and (3) an unconditional $\codim\Sing(V(F)) + 2$ determinantal complexity lower bound for all degree-$n \ge 4$ hypersurfaces.

\subsection{Stage 1: Free Associative Algebra Word-Reversal and Darboux Structure}

Let $\mathcal{A} = \mathbb{C}\langle \mathbf{m}_1,\dots,\mathbf{m}_s \rangle$ be the free associative algebra in $s$ non-commuting generators, and let $\rho : \mathcal{A} \to \mathcal{A}$ be the canonical word-reversal anti-involution defined by $\rho(\mathbf{m}_i) = \mathbf{m}_i$ and $\rho(ab) = \rho(b)\rho(a)$, so that $\rho(\mathbf{m}_{i_1}\cdots\mathbf{m}_{i_k}) = \mathbf{m}_{i_k}\cdots\mathbf{m}_{i_1}$. Since $\rho^2 = \mathrm{id}$ and $\mathrm{char}(\mathbb{C}) \neq 2$, $\mathcal{A} = \mathcal{A}^+ \oplus \mathcal{A}^-$ where $\mathcal{A}^\pm := \{P \in \mathcal{A} : \rho(P) = \pm P\}$ with projections $P^\pm = \frac{1}{2}(P \pm \rho(P))$.

\begin{theorem}[Word-Reversal Identity, Necessary \& Sufficient Conditions, and Darboux Block Structure]\label{thm:wordreversal}
Let $K \in \mathfrak{so}_{d_v}(\mathbb{C})$ be skew-symmetric ($K^T = -K$) of even rank $\rank(K) = 2r \le d_v$, and let $\Phi : \mathcal{A} \to \mathbb{C}^{d_v \times d_v}$ be an algebra homomorphism such that $M_i := \Phi(\mathbf{m}_i)$ satisfies $(M_i K)^T = -M_i K$ for all $1 \le i \le s$. Then:
\begin{enumerate}
  \item \emph{(Universal Word-Reversal Identity):} For every $P \in \mathcal{A}$,
  \[
    \big(\Phi(P) K\big)^T = -\Phi\big(\rho(P)\big) K.
  \]
  \item \emph{(Necessary and Sufficient Annihilation Conditions \& Exact Dimensions):} For any fixed representation $(\Phi, K)$ and any $P = P^+ + P^- \in \mathcal{A}$,
  \[
    y^T \Phi(P) K y = 0 \quad (\forall\, y \in \mathbb{C}^{d_v})
    \;\Longleftrightarrow\;
    \big(\Phi(P) K\big)^T = -\Phi(P) K
    \;\Longleftrightarrow\;
    \Phi(P^-) K = 0.
  \]
  In particular, $P \in \mathcal{A}^+$ ($P^- = 0$) is sufficient for $y^T \Phi(P) K y \equiv 0$ in every admissible representation $(\Phi, K)$. In degree $k \ge 1$, the exact dimensions of the reversal-symmetric subspace $\mathcal{A}_k^+$ (which contains all palindromic words, the special Jordan subalgebra generated under $\{a,b\} = ab+ba$, all single-operator polynomials $p(\mathbf{m}_1)$, and all commutative-scalar pencil powers $(\sum_{i=1}^s \xi_i \mathbf{m}_i)^k$) and the reversal-antisymmetric subspace $\mathcal{A}_k^-$ (containing odd-depth Lie commutators such as $[\mathbf{m}_i, \mathbf{m}_j] = \mathbf{m}_i\mathbf{m}_j - \mathbf{m}_j\mathbf{m}_i$) are
  \[
    \dim_{\mathbb{C}} \mathcal{A}_k^+ = \frac{s^k + s^{\lceil k/2 \rceil}}{2},
    \qquad
    \dim_{\mathbb{C}} \mathcal{A}_k^- = \frac{s^k - s^{\lceil k/2 \rceil}}{2}.
  \]
  \item \emph{(Darboux Block Structure and the $\rank(K) \le 2$ vs.\ $\rank(K) \ge 4$ Dichotomy):} In a Darboux basis where $K = \mathrm{diag}(J_{2r}, 0_{d_v - 2r})$ with $J_{2r} = \left(\begin{smallmatrix} 0 & I_r \\ -I_r & 0 \end{smallmatrix}\right) \in \mathrm{GL}_{2r}(\mathbb{C})$, a matrix $M = \left(\begin{smallmatrix} M_{11} & M_{12} \\ M_{21} & M_{22} \end{smallmatrix}\right) \in \mathbb{C}^{d_v \times d_v}$ partitioned into $(2r) + (d_v - 2r)$ blocks satisfies $(MK)^T = -MK$ if and only if
  \[
    M_{21} = 0 \qquad \text{and} \qquad M_{11} = -Q_{11} J_{2r} \quad \big(Q_{11} \in \mathfrak{so}_{2r}(\mathbb{C})\big).
  \]
  Consequently, the admissible operator space $\mathcal{L}_K := \{M \in \mathbb{C}^{d_v \times d_v} : (MK)^T = -MK\}$ is block upper-triangular of dimension $\dim_{\mathbb{C}} \mathcal{L}_K = \binom{2r}{2} + d_v(d_v - 2r) = d_v^2 - 2r\,d_v + r(2r-1)$, and for any $k$ operators $M^{(1)},\dots,M^{(k)} \in \mathcal{L}_K$,
  \[
    M^{(1)} M^{(2)} \cdots M^{(k)} K = \begin{pmatrix} M_{11}^{(1)} M_{11}^{(2)} \cdots M_{11}^{(k)} J_{2r} & 0 \\ 0 & 0 \end{pmatrix}.
  \]
  Moreover, a sharp rank threshold holds for symmetry breaking:
  \begin{itemize}
    \item If $\rank(K) = 2r \le 2$ (which holds automatically whenever $d_v \in \{1, 2, 3\}$), then $\dim\mathfrak{so}_{2r}(\mathbb{C}) \le 1$, so every $M_{11}^{(t)} = \lambda_t I_{2r}$ is a scalar multiple of $I_{2r}$; hence all $M_{11}^{(t)}$ commute and $\Phi(\mathcal{A}^-) K = 0$ identically (i.e., $y^T \Phi(P) K y \equiv 0$ for \emph{all} $P \in \mathcal{A}$).
    \item Non-trivial commutator escape $\Phi([\mathbf{m}_1, \mathbf{m}_2]) K \neq 0$ exists if and only if $\rank(K) = 2r \ge 4$ (forcing $d_v \ge 4$).
  \end{itemize}
\end{enumerate}
\end{theorem}

\begin{proof}
(1) Since $K^T = -K$, $(M_i K)^T = -M_i K$ is equivalent to $K M_i^T = M_i K$. For any word $\mathbf{w} = \mathbf{m}_{i_1}\cdots\mathbf{m}_{i_k}$, induction on $k$ gives $K M_{i_k}^T \cdots M_{i_1}^T = M_{i_k}\cdots M_{i_1} K = \Phi(\rho(\mathbf{w}))K$, whence $(\Phi(\mathbf{w})K)^T = K^T M_{i_k}^T\cdots M_{i_1}^T = -\Phi(\rho(\mathbf{w}))K$. Linearity extends this to all $P \in \mathcal{A}$.

(2) For $Q = \Phi(P)K$, $y^T Q y = 0$ ($\forall y$) $\iff Q + Q^T = 0 \iff \Phi(P - \rho(P))K = 2\Phi(P^-)K = 0$. Among the $s^k$ words of length $k$, exactly $s^{\lceil k/2 \rceil}$ are palindromes ($\rho(\mathbf{w}) = \mathbf{w}$), while the remaining $s^k - s^{\lceil k/2 \rceil}$ words pair into $\{\mathbf{w}, \rho(\mathbf{w})\}$, contributing one basis element $\mathbf{w} + \rho(\mathbf{w})$ to $\mathcal{A}_k^+$ and one basis element $\mathbf{w} - \rho(\mathbf{w})$ to $\mathcal{A}_k^-$.

(3) With $K = \mathrm{diag}(J_{2r}, 0)$, we have $MK = \left(\begin{smallmatrix} M_{11}J_{2r} & 0 \\ M_{21}J_{2r} & 0 \end{smallmatrix}\right)$. Thus $(MK)^T = -MK$ holds if and only if $(M_{11}J_{2r})^T = -M_{11}J_{2r}$ and $M_{21}J_{2r} = 0$. Since $J_{2r}^{-1} = -J_{2r}$ is invertible, $M_{21} = 0$ and $Q_{11} := M_{11}J_{2r} \in \mathfrak{so}_{2r}(\mathbb{C})$. Closure of block upper-triangular matrices under multiplication yields the product formula and $\dim \mathcal{L}_K = \binom{2r}{2} + 2r(d_v - 2r) + (d_v - 2r)^2$. Finally, if $2r = 0$, $K = 0$; if $2r = 2$, every $Q_{11} \in \mathfrak{so}_2(\mathbb{C})$ has the form $\lambda J_2$, so $M_{11} = -(\lambda J_2)J_2 = \lambda I_2$, which is central in $\mathrm{Mat}_2(\mathbb{C})$ and annihilates all commutators. Conversely, for $2r \ge 4$, taking $Q_{11}^{(1)} = E_{12} - E_{21}$ and $Q_{11}^{(2)} = E_{23} - E_{32}$ in $\mathfrak{so}_{2r}(\mathbb{C})$ gives $[M_{11}^{(1)}, M_{11}^{(2)}] J_{2r} \neq 0$.
\end{proof}

\subsection{Stage 2: Control Theory Hankel Collapse and ABP Width/Hessian Obstructions}

\begin{proposition}[LTI Control Theory Infinite Hankel Collapse]\label{prop:hankel}
In a linear time-invariant system $(M, u, v^T)$ with $v = y$, $u = Ky$ ($K^T = -K$), if the first two Markov parameters vanish identically in $y$ ($h_0 = y^T K y = 0$ and $h_1 = y^T M K y = 0$, i.e., $(MK)^T = -MK$), then every Markov parameter $h_k := v^T M^k u$ ($k \ge 0$), the infinite Hankel matrix $(H_\infty)_{i,j} = h_{i+j}$, and the transfer function $G(s) = v^T(sI - M)^{-1}u$ vanish identically.
\end{proposition}

\begin{proof}
Since $\mathbf{m}_1^k \in \mathcal{A}^+$ for all $k \ge 0$, Theorem~\ref{thm:wordreversal}(2) gives $h_k = y^T M^k K y = 0$ for all $k \ge 0$.
\end{proof}

\begin{theorem}[Palindromic ABP Width Bound and Non-Homogeneous ABP Darboux--Hessian Bound]\label{thm:abp}
Let $F \in \mathbb{C}[x_1,\dots,x_N]_n$ ($n \ge 4$) be computed by a width-$w$ Algebraic Branching Program
\[
  F(X) = v(X)^T A^{(1)}(X) A^{(2)}(X) \cdots A^{(n-2)}(X) u(X)
\]
($v, u \in \mathbb{C}[X]_1^w$, $A^{(t)} \in \mathbb{C}[X]_1^{w \times w}$) satisfying the first-step isotropy conditions $v(X)^T u(X) \equiv 0$ and $v(X)^T A^{(t)}(X) u(X) \equiv 0$ for $1 \le t \le n-2$. Let $W = \Span\{v_i, u_i\}$, $I_W$ be the ideal of $W^\perp = V(I_W)$, and $d_v = \dim\Span\{v_1,\dots,v_w\}$.
\begin{enumerate}
  \item \emph{(Time-Homogeneous and Palindromic ABPs):} If the layer sequence is palindromic ($A^{(t)} = A^{(n-1-t)}$ for all $t$; in particular if $A^{(1)} = \cdots = A^{(n-2)}$) and $\codim_{\mathbb{C}^N}(\Sing_2(V(F))) > \codim_{\mathbb{C}^N}(\Sing(V(F)))$, then
  \[
    w \;\ge\; \codim_{\mathbb{C}^N}\big(\Sing(V(F))\big) + 1.
  \]
  \item \emph{(Non-Homogeneous ABPs with Arbitrary Distinct Layers):} Even when $A^{(1)}, \dots, A^{(n-2)}$ are completely distinct and non-palindromic, at the boundary width $w = \dim W$ (with $d_z := w - d_v$), the leading $I_W^2$ quadratic form on $W^\perp$ depends only on the $2r = \rank(K) \le 2\lfloor d_v/2 \rfloor$ active Darboux coordinates of $y$, so for every $\xi \in W^\perp$:
  \[
    \rank \mathrm{Hess}(F)(\xi) \;\le\; 2r \;=\; \rank(K) \;\le\; 2\lfloor d_v / 2 \rfloor \;\le\; w - d_z.
  \]
  In particular, $w \ge \max_{\xi \in W^\perp} \rank\mathrm{Hess}(F)(\xi) + d_z$, and whenever $\rank(K) \le 2$ (which holds for all $d_v \in \{1, 2, 3\}$), even a non-homogeneous ABP with arbitrary distinct layers forces $\rank\mathrm{Hess}(F)|_{W^\perp} \equiv 0$ ($F \in I_W^3$).
\end{enumerate}
\end{theorem}

\begin{proof}
As in Section~\ref{sec:mainproof}, $v^T u \equiv 0$ and $F \in I_W^2$ yield $\codim\Sing(V(F)) \le \dim W \le w$. At the boundary $w = \dim W$, the $2$-block normal form $v = (y,0)^T$, $u = (Ky,z)^T$ ($K \in \mathfrak{so}_{d_v}(\mathbb{C})$) and $v^T A^{(t)} u \equiv 0$ force $A_{12}^{(t)} \in I_W$ and $A_{11}^{(t,\xi)}(\xi) \in \mathcal{L}_K$ for every layer $t$. Thus $A^{(t)}(X) \pmod{I_W}$ is block lower-triangular, and the leading $I_W^2$ term of $F$ on $W^\perp$ is $y^T A_{11}^{(1,\xi)}(\xi)\cdots A_{11}^{(n-2,\xi)}(\xi) K y$.
If the layers are palindromic, the word lies in $\mathcal{A}^+$, so Theorem~\ref{thm:wordreversal}(2) gives $F \in I_W^3 \implies W^\perp \subseteq \Sing_2(V(F))$, contradicting $\codim\Sing_2(V(F)) > \codim\Sing(V(F)) = w$.
For arbitrary non-palindromic layers, Theorem~\ref{thm:wordreversal}(3) shows that $A_{11}^{(1,\xi)}(\xi)\cdots A_{11}^{(n-2,\xi)}(\xi) K$ is supported exclusively in the top $2r \times 2r$ block, bounding the Hessian rank on $W^\perp$ by $2r \le 2\lfloor d_v/2 \rfloor \le w - d_z$.
\end{proof}

\subsection{Stage 3: Unconditional $\codim(\Sing) + 2$ Determinantal Complexity Bound for Degree $n \ge 4$}

\begin{lemma}[Mandatory $\corank(A_0) = 1$ and Inductive Cross-Term Elimination]\label{lem:generalcorank1}
Let $F \in \mathbb{C}[x_1,\dots,x_N]_n$ ($n \ge 2$) be a homogeneous polynomial satisfying $\codim_{\mathbb{C}^N}(\Sing(V(F))) \ge 5$, and let $F(X) = \det(M(X))$ ($M(X) = A_0 + L(X) \in \mathbb{C}[X]^{m \times m}$) be any affine determinantal representation. Then:
\begin{enumerate}
  \item $\corank(A_0) = m - \rank(A_0) = 1$ (and $w := \rank(A_0) = m - 1$).
  \item In the normal form~\eqref{eq:blockM}--\eqref{eq:schur} with $s_0(X) := b(X)$ and $s_{k+1}(X) := (-1)^{k+1} v(X)^T A(X)^k u(X)$, one has $s_{r-1}(X) \equiv 0$ for all $1 \le r \le n - 1$, and all cross terms in $[t^n]\det(M(tX))$ vanish identically, leaving
  \[
    F(X) = (-1)^{n-1} v(X)^T A(X)^{n-2} u(X).
  \]
\end{enumerate}
\end{lemma}

\begin{proof}
(1) Since $F$ is homogeneous of degree $n \ge 2$, $\det(A_0) = F(0) = 0$, so $r := \corank(A_0) \ge 1$. If $r \ge 2$, then $0 \in M^{-1}(\mathcal{D}_{m-2})$, where $\mathcal{D}_{m-2} = \{B \in \mathbb{C}^{m\times m} : \rank(B) \le m-2\}$ has codimension $(m - (m-2))^2 = 4$. At every $X \in M^{-1}(\mathcal{D}_{m-2})$, all $(m-1)\times(m-1)$ cofactors of $M(X)$ vanish, so $\nabla F(X) = 0$, i.e., $M^{-1}(\mathcal{D}_{m-2}) \subseteq \Sing(V(F))$. Since $0 \in M^{-1}(\mathcal{D}_{m-2})$, every irreducible component $Z$ of $M^{-1}(\mathcal{D}_{m-2})$ containing $0$ has $\codim_{\mathbb{C}^N}(Z) \le 4$, giving $\codim_{\mathbb{C}^N}(\Sing(V(F))) \le 4$, a contradiction. Thus $r = 1$.

(2) Expanding $t^n F(X) = t \big(1 + \sum_{i=1}^w e_i(A) t^i\big)\big(\sum_{j=0}^\infty s_j t^j\big)$ and comparing $[t^r] = 0$ for $1 \le r \le n-1$ gives $s_{r-1} + \sum_{i=1}^{r-1} e_i(A) s_{r-1-i} = 0$, which by induction on $r$ forces $s_0 \equiv s_1 \equiv \cdots \equiv s_{n-2} \equiv 0$. At degree $r = n$, $[t^n]\det(M(tX)) = s_{n-1} + \sum_{i=1}^{n-1} e_i(A) s_{n-1-i} = s_{n-1}$.
\end{proof}

\begin{theorem}[General Degree-$n \ge 4$ Hypersurface Lower Bound]\label{thm:generalplus2}
Let $F \in \mathbb{C}[x_1,\dots,x_N]_n$ be any homogeneous polynomial of degree $n \ge 4$ over $\mathbb{C}$ satisfying
\[
  c_1 := \codim_{\mathbb{C}^N}\big(\Sing(V(F))\big) \ge 5
  \qquad \text{and} \qquad
  \codim_{\mathbb{C}^N}\big(\Sing_2(V(F))\big) > c_1.
\]
Then unconditionally,
\[
  \dc_{\mathbb{C}}(F) \;\ge\; \codim_{\mathbb{C}^N}\big(\Sing(V(F))\big) + 2.
\]
\end{theorem}

\begin{proof}
Suppose $F(X) = \det(A_0 + L(X))$ has size $m$. By Lemma~\ref{lem:generalcorank1}, $\corank(A_0) = 1$, $w = m - 1$, $F(X) = (-1)^{n-1} v(X)^T A(X)^{n-2} u(X)$, and since $n \ge 4$, both $s_1 = -v(X)^T u(X) \equiv 0$ ($[t^2]=0$) and $s_2 = v(X)^T A(X) u(X) \equiv 0$ ($[t^3]=0$) hold identically. Thus $F$ is computed by a width-$w$ time-homogeneous ABP satisfying the hypotheses of Theorem~\ref{thm:abp}(1). By Theorem~\ref{thm:abp}(1), $w = m - 1 \ge c_1 + 1$, and therefore $m \ge c_1 + 2$.
\end{proof}

\begin{thebibliography}{9}

\bibitem{ABV2017}
J.~Alper, T.~Bogart, and M.~Velasco,
\emph{A lower bound for the determinantal complexity of a hypersurface},
Found.\ Comput.\ Math.\ \textbf{17} (2017), no.~3, 829--836.

\bibitem{Atkinson1983}
M.~D.~Atkinson,
\emph{Primitive spaces of matrices of bounded rank.~II},
J.\ Austral.\ Math.\ Soc.\ Ser.~A \textbf{34} (1983), no.~3, 306--315.

\bibitem{EH1988}
D.~Eisenbud and J.~Harris,
\emph{Vector spaces of matrices of low rank},
Adv.\ Math.\ \textbf{70} (1988), no.~2, 135--155.

\bibitem{Grenet2011}
B.~Grenet,
\emph{An upper bound for the permanent versus determinant problem},
Manuscript, \'Ecole Normale Sup\'erieure de Lyon, 2011.

\bibitem{LR2017}
J.~M.~Landsberg and N.~Ressayre,
\emph{Permanent v.\ determinant: an exponential lower bound assuming symmetry and a potential path towards Valiant's conjecture},
Abh.\ Math.\ Semin.\ Univ.\ Hambg.\ \textbf{87} (2017), no.~2, 323--339.

\bibitem{MR2004}
T.~Mignon and N.~Ressayre,
\emph{A quadratic bound for the determinant and permanent problem},
Int.\ Math.\ Res.\ Not.\ IMRN \textbf{2004} (2004), no.~79, 4241--4253.

\bibitem{Sheshadri2026}
K.~Sheshadri,
\emph{The determinantal complexity of the $4 \times 4$ permanent is $15$},
arXiv:2611.06658, 2026 (preprint, not yet peer-reviewed).

\bibitem{Valiant1979}
L.~G.~Valiant,
\emph{Completeness classes in algebra},
Proc.\ 11th Annual ACM Symposium on Theory of Computing (STOC), 1979, pp.~249--261.

\bibitem{Yabe2015}
A.~Yabe,
\emph{Bi-polynomial rank and determinantal complexity},
arXiv:1504.00151, 2015.

\end{thebibliography}

\appendix
\section{Explicit Linear Slice Matrix $M_{L_8}$}\label{app:ml8}

The $8$-dimensional integer subspace $L_8 = \{M_{L_8} t : t \in \mathbb{C}^8\} \subset \mathbb{C}^{16}$ used in Theorem~\ref{thm:sing1} is \emph{normatively defined} by the explicit $16 \times 8$ integer matrix below (which has rank~$8$, and was originally produced in NumPy~2.0.2 via \texttt{np.random.default\_rng(2026).integers(-2, 3, size=(16, 8))}):
\[
  M_{L_8} = \begin{pmatrix}
   2 & -2 & -2 &  1 & -1 &  0 & -2 & -1 \\
   1 & -1 &  2 &  1 &  1 &  2 &  1 & -2 \\
   2 &  1 & -2 & -1 & -2 &  2 &  1 &  2 \\
  -1 &  1 &  1 &  1 & -2 &  0 &  1 &  2 \\
   1 &  0 &  0 & -1 & -2 & -1 & -2 & -1 \\
   2 &  0 & -2 &  0 & -1 &  1 &  2 & -2 \\
   0 &  0 &  2 & -1 & -2 & -2 &  0 &  0 \\
  -1 &  0 &  1 & -1 &  2 & -1 &  2 &  2 \\
   2 &  1 &  0 &  1 & -1 & -1 &  0 &  2 \\
   1 &  0 &  0 &  0 &  0 &  2 & -2 & -1 \\
   0 &  1 &  0 & -1 & -2 & -1 & -2 &  1 \\
   0 & -1 &  0 &  0 &  0 &  0 &  1 & -2 \\
  -2 &  1 &  2 &  0 &  0 &  1 & -2 & -1 \\
  -1 &  0 &  1 &  0 &  2 &  0 & -2 &  1 \\
   1 &  0 &  2 &  0 &  1 & -2 & -2 &  1 \\
   2 &  0 &  1 & -1 &  1 &  0 & -1 & -2
  \end{pmatrix}.
\]

\end{document}
